\chapter{Hyperglycaemia (High Blood Sugar)}

\section*{Definition and Overview}
Hyperglycaemia is the medical term for high blood sugar, a blood glucose level above the normal range. It is the defining feature of diabetes but can also occur temporarily during illness or stress. Persistent hyperglycaemia damages blood vessels and nerves over time, causing complications in the eyes, kidneys, heart, and feet, while very high levels can cause acute emergencies such as diabetic ketoacidosis. Recognising the symptoms early and managing blood sugar well prevents both the short-term and long-term harms.

\section*{Epidemiology}
Around 1.2 million Australians live with diagnosed diabetes, and many more have undiagnosed diabetes or pre-diabetes. High blood sugar is one of the leading causes of cardiovascular disease, kidney failure, and preventable blindness. Type 2 diabetes, the most common form, is increasing with obesity and ageing. Diabetic ketoacidosis affects people with type 1 diabetes and is a leading cause of hospital admission and death in young people with diabetes.

\section*{Aetiology and Pathophysiology}
\begin{itemize}
    \item \textbf{Type 1 diabetes:} the immune system destroys insulin-producing cells, causing absolute insulin deficiency
    \item \textbf{Type 2 diabetes:} insulin resistance and declining insulin production
    \item \textbf{Pre-diabetes:} blood sugar above normal but below diabetic levels
    \item \textbf{Illness and stress:} hormones released during illness raise blood sugar
    \item \textbf{Medications:} steroids and some other drugs raise blood sugar
    \item \textbf{Missed insulin or medicine:} a common cause of acute hyperglycaemia
    \item \textbf{Pathophysiology:} without enough insulin, glucose cannot enter cells and accumulates in the blood
    \item \textbf{Ketoacidosis:} fat breakdown produces ketones and acid, an emergency in type 1 diabetes
\end{itemize}

\section*{Clinical Features}
\begin{itemize}
    \item Increased thirst and frequent urination
    \item Fatigue and weakness
    \item Blurred vision
    \item Dry mouth and headaches
    \item Weight loss in uncontrolled diabetes
    \item Slow healing of cuts and recurrent infections
    \item Ketoacidosis: nausea, vomiting, abdominal pain, fruity breath, rapid breathing, and drowsiness
    \item Many people have no symptoms despite high levels
\end{itemize}

\section*{Differential Diagnosis}
\begin{itemize}
    \item Diabetes mellitus is the main cause
    \item Stress and illness hyperglycaemia, which may be temporary
    \item Medication-induced hyperglycaemia
    \item Endocrine conditions such as Cushing's syndrome
    \item The cause is identified by blood tests and clinical assessment
\end{itemize}

\section*{When to Seek Medical Care}
Anyone with symptoms of high blood sugar or risk factors should be tested, and people with diabetes should follow their sick-day plan during illness.

\textbf{Call 000 (or local emergency services) for:}
\begin{itemize}
    \item Confusion, drowsiness, or loss of consciousness
    \item Rapid deep breathing with a fruity smell (ketoacidosis)
    \item Vomiting with an inability to keep down fluids
    \item Very high blood sugar with dehydration
    \item Chest pain, breathlessness, or signs of stroke
\end{itemize}

\section*{Diagnosis and Investigations}
\begin{itemize}
    \item \textbf{Fasting blood glucose:} 7.0 mmol/L or higher indicates diabetes
    \item \textbf{HbA1c:} reflects average glucose over 3 months; 6.5\% or higher indicates diabetes
    \item \textbf{Oral glucose tolerance test:} for diagnosis where needed
    \item \textbf{Ketone testing:} urine or blood ketones in type 1 diabetes
    \item \textbf{Self-monitoring:} home glucose meters track daily levels
    \item \textbf{Complication screening:} eye, kidney, and foot checks in diabetes
\end{itemize}

\section*{Management}
\begin{itemize}
    \item \textbf{Insulin:} essential in type 1 diabetes, and often in advanced type 2
    \item \textbf{Oral medicines:} metformin and others for type 2 diabetes
    \item \textbf{Healthy eating:} balanced meals with controlled carbohydrates
    \item \textbf{Physical activity:} improves insulin sensitivity
    \item \textbf{Weight management:} significant benefit in type 2 diabetes
    \item \textbf{Sick-day management:} more fluids, more monitoring, and never stopping insulin
    \item \textbf{Blood sugar monitoring:} regular testing and review
    \item \textbf{Education and support:} diabetes education improves outcomes
\end{itemize}

\section*{Complications}
\begin{itemize}
    \item Diabetic ketoacidosis and hyperosmolar hyperglycaemic state
    \item Eye disease, kidney disease, and nerve damage
    \item Heart attack, stroke, and peripheral artery disease
    \item Foot ulcers and amputation
    \item Infections and poor wound healing
    \item Reduced quality of life
\end{itemize}

\section*{Prognosis and Outcomes}
With good control, people with diabetes live long, healthy lives. Every 1\% reduction in HbA1c reduces the risk of complications by around a third, and early treatment of pre-diabetes can prevent progression. Modern technology, including continuous glucose monitors, has transformed management. The keys to good outcomes are early detection, consistent treatment, and regular complication screening.

\section*{Prevention and Self-Care}
\begin{itemize}
    \item Maintain a healthy weight and be physically active
    \item Eat a balanced diet and limit sugar and refined carbohydrates
    \item Get tested if you have risk factors
    \item Take diabetes medicines and insulin as prescribed
    \item Monitor blood sugar and keep a record
    \item Attend annual eye, kidney, and foot screening
    \item Have a sick-day plan and know the signs of ketoacidosis
    \item Seek help early for high readings or symptoms
\end{itemize}

\section*{Sources and Further Reading}
The following reputable sources provide further information for healthcare students and patients seeking reliable, current advice about hyperglycaemia.

\begin{itemize}
    \item Diabetes Australia. \textit{High blood sugar information.} https://www.diabetesaustralia.com.au/
    \item Healthdirect. \textit{Hyperglycaemia (high blood sugar).} https://www.healthdirect.gov.au/high-blood-sugar
    \item National Diabetes Services Scheme. \textit{Diabetes support and products.} https://www.ndss.com.au/
    \item NHS. \textit{Hyperglycaemia.} https://www.nhs.uk/conditions/high-blood-sugar-hyperglycaemia/
    \item World Health Organization. \textit{Diabetes.} https://www.who.int/
    \item Mayo Clinic. \textit{Hyperglycemia in diabetes.} https://www.mayoclinic.org/
\end{itemize}
